\documentclass[11pt]{article}
\usepackage[margin=1in]{geometry}
\usepackage{amsmath,amssymb,amsthm}
\usepackage{enumitem}
\usepackage{xcolor}
\usepackage{parskip}
\usepackage{booktabs}
\usepackage{hyperref}
\hypersetup{colorlinks=true, urlcolor=blue!45!black, linkcolor=blue!45!black}
\title{\vspace{-2em}\textbf{Statistical Tests for Hypotheses H1 \& H3}\\[0.3em]\large Ordered-trend tests, rank association, and donor-level inference\\[0.6em]\normalsize\textit{Genomics Hackathon --- methods primer}}
\author{}\date{}

\newcommand{\dez}{\text{dez}}
\newcommand{\plast}{\text{plast}}

\begin{document}
\maketitle
\tableofcontents

\bigskip

\noindent\fbox{\parbox{0.97\textwidth}{%
\textbf{Trend across ordered disease stages, and association between two per-cell scores --- with the donor as the unit of inference.}\\[0.4em]
A teaching companion to the de-zonation analysis.\\
\textbf{H1:} does the de-zonation score change monotonically across ordered stages?\\
\textbf{H3:} do de-zonation and plasticity go together, after controlling for stage?\\[0.3em]
Cohort: $\approx 47$ donors $\cdot$ hepatocytes $\cdot$ ordered disease stages.\\
\textcolor{blue!45!black}{\textbf{Golden rule:} the unit of inference is the \emph{donor}, never the cell.}}}

\section{The one rule that governs everything: the donor is the unit}

Before any test, fix the most important idea. We have $\approx 47$ donors. Each donor contributes thousands of single cells, and from those cells we compute \emph{one summary per donor} --- a de-zonation score. The biological question (``does de-zonation track disease?'') is a statement about \emph{donors}, so the sample size that matters is
\[
n \approx 47, \qquad \text{not} \qquad n \sim 5\times 10^{5}\ \text{cells}.
\]

\paragraph{Pitfall --- pseudoreplication.} If you run a test on cell-level values (treating each cell as an independent observation), you inflate $n$ from $\approx 47$ to $\approx 500{,}000$. Cells from the same donor are \emph{not} independent --- they share that donor's genetics, treatment, batch and biology. Concretely: a real per-donor effect that an honest donor-level test rates at, say, $p\approx 0.2$ can come out as $p<10^{-40}$ once you wrongly treat half a million cells as independent --- that tiny number reflects how many \emph{cells} you sequenced, not how many \emph{patients} actually support the claim. This error is called \textbf{pseudoreplication}, and it is the single easiest way to fool yourself in single-cell work. Every test in this primer is run on \emph{one number per donor}.

\paragraph{The two-step pattern.}
\begin{enumerate}[label=(\arabic*),leftmargin=*]
  \item Collapse cells $\rightarrow$ one score per donor.
  \item Run the donor-level test below.
\end{enumerate}
Cell-level structure (like score-vs-score correlation \emph{within} a donor) is the exception in Part~B, and even there we aggregate back up to the donor.

With that nailed down, H1 asks about a \emph{trend} over an ordered axis, and H3 asks about an \emph{association} between two scores. Different shapes of question, different tools.

\section{Part A --- H1: does the score change monotonically across ordered stages?}

\paragraph{Setup.} Each donor has (a) a stage label that is \emph{ordered} (e.g.\ healthy $\rightarrow$ early $\rightarrow$ mid $\rightarrow$ advanced) and (b) one de-zonation score. We are not asking ``are the groups different?'' but the sharper ``does the score \emph{rise (or fall) in step} with stage?'' Ordering is information --- tests that use it are more powerful than tests that ignore it.

\paragraph{Where the per-donor score comes from (and why the gene set matters here).} The de-zonation score is read from a \emph{gene program}, not a few markers: per cell, $\text{coord} = \operatorname{mean}_z(\text{pericentral genes}) - \operatorname{mean}_z(\text{periportal genes})$, aggregated to one number per donor. We run every test below \textbf{twice} --- on the transcriptome-wide \texttt{full} set ($\approx$1{,}273 PC $/$ 364 PP genes; the primary, because averaging hundreds of genes makes the ruler robust to any single disease-corrupted marker) and on the \texttt{paper2\_landmark} set (20 PC $/$ 20 PP, Paper 2's exact hepatocyte landmark genes; the interpretable audit). Agreement of the two is itself evidence the trend is real, not a gene-set artefact. Two consequences for the statistics: (i)~the PC/PP count imbalance does \emph{not} bias the score, because each arm is a \emph{mean} (not a sum) of z-scores --- though the smaller PP arm is noisier, so we standardise each arm to unit variance before subtracting; (ii)~for the H2 ``which genes lose their zonal pattern'' test the coordinate is built from these same genes, so that test must use a \emph{held-out} random gene split repeated $K\!\approx\!20$--$50\times$ --- which is only meaningful with the \texttt{full} set, not 20$+$20 landmarks.

\subsection{Spearman rank correlation $\rho$}

\textit{The question.} Each donor gives us a pair: a disease stage and a de-zonation score. We want one number that answers ``as stage goes up, does the score climb with it?'' --- and we want it to be robust, i.e.\ to \emph{not} assume the climb is a straight line or that the stages are evenly spaced. We build that number in three steps: start from the ordinary Pearson correlation, replace the values by their ranks, then simplify.

\paragraph{Step 1 --- start from Pearson, and see why it's wrong here.} Pearson's $r$ measures how much two lists move \emph{together}, rescaled to $[-1,1]$:
\[
r = \frac{\sum_i (x_i-\bar x)(y_i-\bar y)}{\sqrt{\sum_i (x_i-\bar x)^2}\;\sqrt{\sum_i (y_i-\bar y)^2}}
\;=\; \frac{\text{how much } x \text{ and } y \text{ move together}}{(\text{spread of }x)\times(\text{spread of }y)}.
\]
Each numerator term $(x_i-\bar x)(y_i-\bar y)$ is \emph{positive} when $x_i$ and $y_i$ land on the same side of their means (both above, or both below) and \emph{negative} when they oppose; summed, the numerator is large when the two rise and fall together. The denominator just divides out the units, so rescaling $y$ from metres to centimetres leaves $r$ unchanged. Two features make Pearson the wrong tool for ``does the score climb with \emph{stage}?'':
\begin{itemize}[leftmargin=*,topsep=2pt]
  \item \textbf{It assumes a straight line.} The curve $y=(1,4,9,16)$ rises flawlessly with $x=(1,2,3,4)$, yet Pearson scores it \emph{below} $1$ simply because the rise isn't linear.
  \item \textbf{It trusts the stage \emph{numbers}, not just their order.} Stages are an ordering --- stage 4 is later than stage 2, not ``twice'' it. But relabelling the same four ordered stages from $(1,2,3,4)$ to $(1,2,10,100)$ changes Pearson's answer, because it reads the labels as real distances.
\end{itemize}

\paragraph{Step 2 --- replace each value by its rank.} A value's \emph{rank} is just its position after sorting: smallest $\to 1$, next $\to 2$, and so on. For $x=(50,20,80,40)$ the sorted order is $20<40<50<80$, so $\operatorname{rank}(x)=(3,1,4,2)$. Spearman's $\rho$ is simply Pearson's $r$ run on the ranks:
\[
\rho \;=\; r\big(\operatorname{rank}(x),\,\operatorname{rank}(y)\big).
\]
Ranks keep only the order and discard the magnitudes, which fixes both problems at once. The curve $y=(1,4,9,16)$ has ranks $(1,2,3,4)$ identical to $x$'s, so $\rho=1$ --- a perfect \emph{monotone} trend. And $(1,2,3,4)$ and $(1,2,10,100)$ share the same ranks, so arbitrary stage spacing stops mattering; a lone outlier can move a rank by only one position, so it can't dominate either.

\paragraph{Step 3 --- the textbook shortcut (assumes no ties).} When no two values tie, each rank list is simply a permutation of $1,2,\dots,n$, so both rank-means and both rank-standard-deviations are fixed known numbers. Substituting those constants into the Pearson formula of Step 1 and simplifying collapses the whole thing into a count of how far the two rankings disagree:
\[
\rho = 1 - \frac{6\sum_i d_i^2}{n\,(n^2-1)},\qquad d_i = \operatorname{rank}(x_i)-\operatorname{rank}(y_i).
\]
Read it directly: $d_i$ is how many places donor $i$ moves between the two orderings. If the orderings agree perfectly, every $d_i=0$, the fraction is $0$, and $\rho=1$. If they are exactly reversed, $\sum d_i^2$ hits its maximum, the fraction reaches $2$, and $\rho=-1$. The constants $6$ and $n(n^2-1)$ are precisely what is needed to map ``total disagreement'' onto the $[-1,1]$ scale. This is an algebraic identity --- it equals Step 2 \emph{whenever} there are no ties, and it is the famous form people quote.

\begin{itemize}[leftmargin=*]
  \item \textbf{Sign:} $\rho > 0 \Rightarrow$ score tends to rise with stage; $\rho < 0 \Rightarrow$ it falls.
  \item \textbf{Magnitude:} $|\rho| = 1$ is a perfect monotone relation (ranks agree exactly); $\rho = 0$ means no monotone trend.
\end{itemize}

\paragraph{Is it significant? --- where the $t$-statistic comes from.} The cutoff is not pulled from nowhere: it is the ordinary Pearson significance test, which is secretly the test of whether the slope of a simple regression is zero. Work with centred variables $X_i=x_i-\bar x$, $Y_i=y_i-\bar y$ and fit the line $Y_i=\beta X_i+\varepsilon_i$. Testing $H_0:\rho=0$ is then the same as testing $H_0:\beta=0$ --- a zero best-fit slope means $Y$ has no tendency to rise or fall with $X$. The least-squares slope can be written in terms of $r$ itself:
\[
\hat\beta=\frac{\sum_i X_iY_i}{\sum_i X_i^2}=r\sqrt{\frac{\sum_i Y_i^2}{\sum_i X_i^2}}.
\]
We judge that slope against its own noise, $t=\hat\beta/\operatorname{SE}(\hat\beta)$, with $\operatorname{SE}(\hat\beta)=s/\sqrt{\sum_i X_i^2}$ and residual variance $s^2=\mathrm{SSE}/(n-2)$. The $n-2$ is the degrees of freedom: two of the $n$ data points were spent estimating an intercept and a slope. The single identity that makes everything collapse is
\[
\mathrm{SSE}=\sum_i Y_i^2\,(1-r^2),
\]
i.e.\ the line explains a fraction $r^2$ of $Y$'s variation and leaves $1-r^2$ as error. Substituting $s$ and $\hat\beta$ into $t$, the clutter $\sum_i X_i^2$ and $\sum_i Y_i^2$ cancels, leaving
\[
t=\frac{\hat\beta}{\operatorname{SE}(\hat\beta)}=r\sqrt{\frac{n-2}{1-r^2}}.
\]
Read it as $t=\dfrac{\text{observed correlation}}{\text{noise a null correlation would show}}$: the $\sqrt{n-2}$ says a given $r$ is more convincing with more donors, and the $1/\sqrt{1-r^2}$ says the evidence explodes as $r\to\pm1$, because a near-perfect trend almost never arises by chance. Under normal-linear assumptions $t\sim t_{\,n-2}$, and $p=\Pr\!\big(|T_{n-2}|\ge|t_{\text{obs}}|\big)$.

\paragraph{\dots but for Spearman this is only an approximation.} Spearman is $\rho=r(\operatorname{rank}(x),\operatorname{rank}(y))$, so the formula above is simply reused with $\rho$ in place of $r$. The catch: ranks are \emph{not} normally distributed, so the $t_{\,n-2}$ result is only a large-$n$ approximation. The \emph{exact} null is combinatorial --- if stage and score are unrelated, every pairing of the two rank lists is equally likely. That is precisely the permutation test: hold one rank list fixed, reshuffle the other many times, recompute $\rho$, and see how often chance alone matches or beats the observed value (the label-shuffle control in \S\ref{sec:perm}). For our modest, heavily-tied $n$ we prefer this exact $p$-value over the $t$ approximation.

\paragraph{Ties --- and we have many.} The Step-3 shortcut assumes no ties, but our stages are \emph{heavily} tied: dozens of donors share each stage label. With ties, the ranks are no longer a clean permutation of $1,\dots,n$, so the $6\sum d_i^2$ form is simply wrong. The fix is to give tied items their shared \emph{average} rank (``mid-ranks'') and fall back to the Step-2 definition --- ordinary Pearson computed on those mid-ranks. \texttt{scipy.stats.spearmanr} does exactly this automatically, which is why in code we call it rather than typing the $6\sum d_i^2$ shortcut by hand.

\paragraph{When to prefer Spearman over Pearson.} Prefer Spearman when the $x$-axis is ordinal (stages), when the relationship may be monotone-but-curved, when there are outliers, or when scores are not normally distributed --- all true here. Pearson would impose linearity and equal spacing on stages we have no right to assume.

\subsubsection*{Worked mini-example (6 donors)}

Suppose six donors at stages $1$--$6$ (already ordered) have de-zonation scores. To make the ranks honest, sort the scores: $0.12, 0.18, 0.20, 0.29, 0.31, 0.40 \rightarrow$ their owners are A, C, B, E, D, F, giving score-ranks $\text{A}=1,\ \text{B}=3,\ \text{C}=2,\ \text{D}=5,\ \text{E}=4,\ \text{F}=6$.

\begin{center}
\begin{tabular}{lcccrr}
\toprule
Donor & Stage (rank $x$) & Score & Score rank ($y$) & $d = x-y$ & $d^2$ \\
\midrule
A & 1 & 0.12 & 1 & $0$  & 0 \\
B & 2 & 0.20 & 3 & $-1$ & 1 \\
C & 3 & 0.18 & 2 & $+1$ & 1 \\
D & 4 & 0.31 & 5 & $-1$ & 1 \\
E & 5 & 0.29 & 4 & $+1$ & 1 \\
F & 6 & 0.40 & 6 & $0$  & 0 \\
\midrule
\multicolumn{5}{r}{$\sum d_i^2 =$} & 4 \\
\bottomrule
\end{tabular}
\end{center}

Now:
\begin{align*}
\sum_i d_i^2 &= (1-1)^2+(2-3)^2+(3-2)^2+(4-5)^2+(5-4)^2+(6-6)^2 \\
             &= 0+1+1+1+1+0 = 4, \\[0.4em]
\rho &= 1 - \frac{6\cdot 4}{6\,(36-1)} = 1 - \frac{24}{210} = 1 - 0.114 = 0.886.
\end{align*}

A strong positive monotone trend: scores rise with stage, with two minor local swaps (B/C and D/E) costing us a little from a perfect $1.0$.

\subsection{Jonckheere--Terpstra test (the dedicated ordered-alternative test)}

\textit{Intuition.} Spearman treats stage as a single ranked variable. The Jonckheere--Terpstra (JT) test is purpose-built for the case where donors fall into $k \ge 3$ \emph{ordered} groups (the stages) and you want to test the directional alternative ``medians don't decrease as you walk along the order.'' It is the ordered cousin of Kruskal--Wallis.

\paragraph{Hypotheses.}
\begin{align*}
H_0&:\ \text{all } k \text{ stage distributions are identical.}\\
H_1&:\ M_1 \le M_2 \le \dots \le M_k \quad \text{(non-decreasing, with at least one strict increase).}
\end{align*}

\paragraph{How the $J$ statistic is built.} For every ordered pair of stages $(a < b)$, count how often a value from the later stage exceeds a value from the earlier stage --- this is exactly the Mann--Whitney $U$ count for that pair (ties count $\tfrac12$). Sum these counts over all $k(k-1)/2$ ordered pairs:
\[
J = \sum_{a<b} U_{ab}, \qquad U_{ab} = \#\{\, x_b > x_a : x_a \in \text{stage } a,\ x_b \in \text{stage } b \,\}.
\]
If the score systematically climbs with stage, later-stage values beat earlier-stage values most of the time, so every $U_{ab}$ is large and $J$ is large.

\subsubsection*{Worked mini-example (3 stages, 2 donors each)}
Scores by stage: stage~1 $=\{0.1,0.3\}$, stage~2 $=\{0.2,0.5\}$, stage~3 $=\{0.4,0.6\}$. For each ordered stage pair, count how many later-stage values beat earlier-stage values:
\begin{itemize}[leftmargin=*,topsep=2pt]
  \item $U_{12}$: $\{0.2,0.5\}$ vs $\{0.1,0.3\}$ --- $0.2{>}0.1$, $0.5{>}0.1$, $0.5{>}0.3$ win, $0.2{<}0.3$ loses $\Rightarrow 3$.
  \item $U_{13}$: $\{0.4,0.6\}$ vs $\{0.1,0.3\}$ --- all four win $\Rightarrow 4$.
  \item $U_{23}$: $\{0.4,0.6\}$ vs $\{0.2,0.5\}$ --- $0.4{>}0.2$, $0.6{>}0.2$, $0.6{>}0.5$ win, $0.4{<}0.5$ loses $\Rightarrow 3$.
\end{itemize}
So $J = 3+4+3 = 10$, against the maximum of $3\text{ pairs}\times 4 = 12$: nearly every later cell outranks every earlier one, a strong upward march. Under the null we will see this is well above the chance value $\mu_J = 6$.

\paragraph{Why JT beats Kruskal--Wallis here.} Kruskal--Wallis only asks ``are the groups different \emph{somewhere}?'' --- it spreads its power across all possible patterns of difference. JT spends all its power on the one pattern we predicted in advance (a march in stage order). When that a-priori ordering is real, JT detects it with a smaller sample. The cost: if the true effect is non-monotone (e.g.\ up then down), JT can miss it while KW would not.

\paragraph{Significance --- and where $\mu_J$ comes from.} Under $H_0$, $J$ has a known mean and variance. The mean is not mysterious. The total number of \emph{cross-stage} pairs is every pair minus the within-stage ones, halved:
\[
M = \frac{N^2 - \sum_s n_s^2}{2}.
\]
Under $H_0$ the two stages in any such pair are interchangeable, so ``later beats earlier'' is a fair coin --- probability $\tfrac12$. The expected number of wins is therefore half the pairs:
\[
\mu_J = \mathbb{E}[J] = \tfrac12 M = \frac{N^2 - \sum_s n_s^2}{4}.
\]
(In the worked example $M = (36-12)/2 = 12$ cross-pairs, so $\mu_J = 6$ --- exactly the chance value our $J=10$ beat.) The variance is the messier bookkeeping of how much that sum of (correlated) coin-flips can wobble:
\[
\sigma_J^2 = \operatorname{Var}(J) = \frac{N^2(2N+3) - \sum_{s} n_s^2(2n_s+3)}{72}.
\]
A normal approximation then standardises $J$ into a $z$-score,
\[
z = \frac{J - \mu_J}{\sigma_J},
\]
and gives a one-sided $p$-value. For small or heavily-tied samples (ours) the normal approximation is rough, so prefer an exact or permutation null --- shuffle stage labels, recompute $J$ thousands of times. Use a \emph{one-sided} test: we predicted the direction in advance.

\subsection{Donor-level bootstrap confidence interval on the trend}

\textit{Intuition.} A $p$-value says ``is the trend distinguishable from zero?'' A confidence interval says ``how big is it, and how sure are we?'' The bootstrap builds a CI by simulating the sampling variability directly: pretend our $47$ donors are the population, draw new pseudo-cohorts from them, and watch how much $\rho$ wobbles.

\begin{quote}
\ttfamily
\textbf{BOOTSTRAP CI for the trend statistic (rho)}\\
inputs : donors = [(stage\_i, score\_i)] for i in 1..n \ \ \# n $\sim$ 47\\
B = 10000 \ \ \# resamples\\
for b in 1..B:\\
\hspace*{2em}sample n donors WITH REPLACEMENT from `donors' \ \ \# resample DONORS\\
\hspace*{2em}rho\_b = spearman(stage, score) on the resampled set\\
\hspace*{2em}store rho\_b\\
CI = ( percentile(rho\_b, 2.5) , percentile(rho\_b, 97.5) ) \ \ \# 95\% interval\\
report rho\_observed and CI
\end{quote}

\paragraph{Why the percentiles give a CI.} The $B=10{,}000$ resampled $\rho_b$ values trace out the \emph{sampling distribution} of $\rho$ --- the spread of answers we would get from cohorts like ours. The middle $95\%$ of that spread is bracketed by its $2.5$th and $97.5$th percentiles, so those two numbers \emph{are} the $95\%$ CI. For example, if the resamples cluster around $0.5$ with $2.5$th/$97.5$th percentiles at $0.28$ and $0.69$, we report ``$\rho = 0.50$, $95\%$ CI $[0.28,\,0.69]$.'' Because the interval excludes $0$, the trend is significant at the $5\%$ level in the bootstrap sense; the width $0.28$--$0.69$ communicates how precise (or not) the estimate is --- something a bare $p$-value never tells you.

\paragraph{Resample donors, not cells.} The resampling unit must match the inference unit. If you bootstrapped cells, you would model only within-donor noise and again commit pseudoreplication --- the CI would be far too narrow. Resampling whole donors captures donor-to-donor variability, which is the variability our conclusion actually depends on.

\subsection{Permutation / label-shuffle negative control}\label{sec:perm}

\textit{Intuition.} How large would $\rho$ (or $J$) look by chance if stage and score had nothing to do with each other? Break the link on purpose: keep the scores, randomly reassign the stage labels across donors, and recompute the statistic. Repeat thousands of times to build the empirical null distribution --- the spread of values produced by pure chance.

\begin{quote}
\ttfamily
\textbf{PERMUTATION p-value (label-shuffle null)}\\
observed = spearman(stage, score) \ \ \# real labels\\
null = []\\
for p in 1..P: \ \ \# P = 10000\\
\hspace*{2em}stage\_shuffled = random permutation of stage labels across donors\\
\hspace*{2em}null.append( spearman(stage\_shuffled, score) )
\end{quote}

The two-sided permutation $p$-value (with the standard ``$+1$'' to avoid $p=0$) is
\[
p = \frac{\#\{\, |\rho^*| \ge |\rho| \,\} + 1}{B + 1},
\]
where $\rho^*$ ranges over the $B$ shuffled statistics and $\rho$ is the observed value.

\paragraph{What a valid negative control proves.} Two things. (1) It gives an \emph{assumption-free} $p$-value --- no normality, no large-$n$ approximation, just the data's own labels. (2) Run on a quantity that should be null (e.g.\ shuffled labels, or a score known to be unrelated), it confirms your pipeline does not manufacture signal from noise: the observed-on-shuffled statistic should land in the bulk of the null, $p \approx 0.5$. A pipeline that ``finds'' a trend in shuffled data is broken.

\section{Part B --- H3: do de-zonation and plasticity go together, controlling for stage?}

\paragraph{Setup.} Now \emph{two} per-cell scores: de-zonation and plasticity. H3 asks whether they are associated. The trap: both scores tend to rise with disease stage, so a naive \emph{pooled} correlation will look positive even if they are unrelated within any single donor. The whole job of Part~B is to remove that confound.

\subsection{Mann--Whitney $U$ / Wilcoxon rank-sum test}

\textit{Intuition.} Split cells (or donors) into two groups --- say de-zonated vs zonated --- and ask ``do plasticity values in one group tend to be larger?'' Pool all values, rank them from low to high, and check whether one group hogs the high ranks. No means, no normality --- just ``who tends to outrank whom.''

\paragraph{The $U$ statistic.} For groups of size $n_1, n_2$, count every pair (one from each group) where the group-1 value exceeds the group-2 value (ties $= \tfrac12$): that count is $U_1$. There are $n_1 n_2$ such cross-group pairs in total, and each is won by exactly one group, so $U_1 + U_2 = n_1 n_2$.

\emph{Why the rank-sum shortcut works.} Computing $U_1$ from pairwise comparisons is tedious; the rank-sum gives it for free. Pool both groups and rank everything; let $R_1$ be the sum of group~1's ranks. If group~1 held the very \emph{smallest} values, it would occupy ranks $1,2,\dots,n_1$, whose sum is the minimum possible, $\tfrac{n_1(n_1+1)}{2}$ --- and in that case group~1 wins \emph{zero} cross-pairs, $U_1=0$. Every rank group~1 carries \emph{above} that floor corresponds to exactly one group-2 value it has overtaken. Hence
\[
U_1 = R_1 - \frac{n_1(n_1+1)}{2}, \qquad U_1 + U_2 = n_1 n_2.
\]

\begin{itemize}[leftmargin=*]
  \item No normality assumption --- it works on ranks, robust to skew and outliers.
  \item \textbf{Ties:} use mid-ranks and a tie-corrected variance for the normal approximation.
  \item \textbf{Relation to AUC:}
  \[
  \text{AUC} = \frac{U_1}{n_1 n_2}
  \]
  --- the probability a random group-1 value outranks a random group-2 value. This is the ROC area exactly.
  \item \textbf{Effect size (rank-biserial):}
  \[
  r = 1 - \frac{2U}{n_1 n_2} \;=\; 2\cdot\text{AUC} - 1 \;=\; \frac{U_1 - U_2}{n_1 n_2},
  \]
  ranging $-1\dots+1$. Always report it beside the $p$-value.
\end{itemize}

\subsubsection*{Worked mini-example}

Plasticity in de-zonated donors: $\{7, 9, 12\}$; in zonated donors: $\{4, 6, 8\}$. Pool and rank:

\begin{center}
\begin{tabular}{lcc}
\toprule
Value & Rank & Group \\
\midrule
4  & 1 & zonated \\
6  & 2 & zonated \\
7  & 3 & de-zonated \\
8  & 4 & zonated \\
9  & 5 & de-zonated \\
12 & 6 & de-zonated \\
\bottomrule
\end{tabular}
\end{center}

Rank-sum of the de-zonated group $R_1 = 3+5+6 = 14$. Then
\begin{align*}
U_1 &= R_1 - \frac{n_1(n_1+1)}{2} = 14 - \frac{3\cdot 4}{2} = 14 - 6 = 8,
      & U_2 &= n_1 n_2 - U_1 = 3\cdot 3 - 8 = 1, \\[0.4em]
\text{AUC} &= \frac{U_1}{n_1 n_2} = \frac{8}{9} = 0.89,
      & r &= 1 - \frac{2U_2}{n_1 n_2} = 2\cdot 0.89 - 1 = +0.78.
\end{align*}

De-zonated donors almost always out-rank zonated ones on plasticity --- a large positive effect (with $n=3$ per group the $p$-value is, of course, not significant; effect size and $p$ tell different stories).

\subsection{Within-stratum / within-donor correlation}

\textit{Intuition --- the confound, drawn out.} Imagine plasticity and de-zonation are \emph{independent within every donor}, but both drift upward as disease advances. Pool all cells across donors and you get a strong positive correlation --- entirely from the shared stage trend, a textbook Simpson's paradox / confounding.

\emph{A concrete picture.} Take an early-stage donor whose cells scatter around $(\dez,\plast)\approx(0.2,0.2)$ with no internal pattern, and a late-stage donor scattering around $(0.8,0.8)$, again with no internal pattern. Within each donor the correlation is $\approx 0$ --- the two scores are unrelated. But pooled, the data form two clouds sitting on the diagonal, and a single correlation through both reads $\rho\approx +0.9$. That ``$0.9$'' is pure stage confound: it says only that sicker donors have higher \emph{everything}, not that de-zonation and plasticity track each other within a patient. The fix: \emph{never pool}. Measure the correlation \emph{inside} each donor (or each stage stratum), where stage is held fixed, then aggregate those clean within-donor numbers.

\begin{quote}
\ttfamily
\textbf{WITHIN-DONOR correlation, then aggregate}\\
for each donor d:\\
\hspace*{2em}cells\_d = all cells of donor d \ \ \# stage is constant here\\
\hspace*{2em}rho\_d = spearman(dezonation, plasticity) on cells\_d\\
summary:\\
\hspace*{2em}mean\_rho = mean(rho\_d over donors) \ \ \# average within-donor assoc.\\
\hspace*{2em}frac\_pos = fraction of donors with rho\_d > 0 \ \ \# consistency / direction\\
\hspace*{2em}test = sign test or Wilcoxon signed-rank on \{rho\_d\} vs 0 \ \ \# n $\sim$ 47
\end{quote}

\paragraph{Why this works.} Within one donor, stage is constant --- it cannot drive a spurious correlation. Each donor gives one honest $\rho_d$. We then have $\approx 47$ such numbers and test whether their center sits above zero. Reporting both mean $\rho$ (strength) and \% donors with $\rho_d>0$ (consistency) is far more convincing than a single pooled coefficient.

\subsection{Linear model with covariates (OLS)}

\textit{Intuition.} The within-donor recipe above is the ``by hand'' version of regression. A single linear model can absorb the stage and donor effects as covariates and report the de-zonation--plasticity link that \emph{remains} after they are removed:
\[
\plast \sim \dez + C(\text{stage}) + C(\text{donor}).
\]

\begin{itemize}[leftmargin=*]
  \item $C(\text{stage})$ and $C(\text{donor})$ are categorical (dummy) terms --- they soak up each stage's and each donor's baseline plasticity.
  \item The coefficient $\beta$ on $\dez$ is the association \emph{after} removing stage and donor effects: ``holding stage and donor fixed, a one-unit rise in de-zonation moves plasticity by $\beta$.''
  \item Its $p$-value (and CI) test $H_0:\beta = 0$ --- no partial association.
\end{itemize}

\paragraph{What the dummy terms actually do.} ``$C(\text{donor})$'' expands into one $0/1$ indicator column per donor, which lets the model fit a \emph{separate intercept} for every donor --- its own baseline plasticity. With those baselines absorbed, the single shared slope $\beta$ can only be estimated from how $\plast$ and $\dez$ move \emph{together within each donor's own cloud}, because the between-donor differences have been subtracted away. That is literally the within-donor correlation of the previous section, now pooled into one estimate with one $p$-value and one CI. $C(\text{stage})$ does the same job for stage.

\paragraph{Caveats.} OLS assumes a roughly linear, additive relationship and well-behaved residuals. And with cell-level rows you still owe an honest treatment of the donor clustering --- the $C(\text{donor})$ term, or a mixed-effects model with a donor random effect, or donor-clustered standard errors --- otherwise the $\dez$ $p$-value is over-confident: pseudoreplication in regression clothing.

\subsection{Shared section --- rank-based vs parametric, multiple testing, effect sizes}

\begin{center}
\begin{tabular}{lll}
\toprule
Aspect & Rank-based (Spearman, MWU, JT) & Parametric (Pearson, $t$-test, OLS) \\
\midrule
Assumes normality?            & No               & Yes (residuals) \\
Robust to outliers?           & Yes              & No \\
Uses exact spacing of values? & No (only order)  & Yes \\
Handles covariates?           & Awkwardly        & Naturally (regression) \\
Power when assumptions hold   & Slightly lower   & Highest \\
\bottomrule
\end{tabular}
\end{center}

Rule of thumb here: default to rank-based for the headline tests (ordinal stages, skewed scores, small $n$), and use OLS/mixed models when you need to adjust for several covariates at once. Report both when feasible --- agreement is reassuring.

\paragraph{Multiple testing --- BH-FDR.} If you test many signatures / many scores, the chance of at least one false positive explodes: at $\alpha=0.05$, even $20$ truly-null tests give a coin-flip chance of a spurious ``hit.'' Control the false discovery rate (the expected fraction of your declared hits that are false) with Benjamini--Hochberg: sort the $m$ $p$-values $p_{(1)} \le \dots \le p_{(m)}$, find the largest $k$ with
\[
p_{(k)} \le \frac{k}{m}\,q \qquad (\text{e.g.\ } q = 0.05),
\]
and declare tests $1\dots k$ significant.

\emph{Why the sloped threshold.} If all $m$ nulls were true, the sorted $p$-values would lie roughly on the straight line $k/m$. BH compares each $p_{(k)}$ to that line scaled by $q$, and keeps everything up to the last point still under it --- letting genuinely small $p$-values ``earn'' a more lenient cutoff than the flat Bonferroni level $q/m$.

\emph{Worked example} ($m=5$, $q=0.05$). Sorted $p$-values $0.001,\,0.013,\,0.025,\,0.030,\,0.45$; thresholds $\tfrac{k}{m}q = 0.01,\,0.02,\,0.03,\,0.04,\,0.05$. Checking: $0.001{<}0.01$, $0.013{<}0.02$, $0.025{<}0.03$, $0.030{<}0.04$ all pass, $0.45{<}0.05$ fails. The largest passing $k$ is $4$, so we reject tests $1$--$4$. Bonferroni would have used the single cutoff $0.05/5 = 0.01$ and rejected only the first --- BH recovers three extra real hits while still controlling the false-discovery rate.

\paragraph{Always report effect sizes beside $p$-values.} A $p$-value answers ``could this be chance?'' --- not ``is it big enough to matter?'' With $\approx 47$ donors a tiny, biologically trivial effect can be significant, and a real effect can miss significance. So pair every $p$ with its effect size: $\rho$ for Spearman, rank-biserial $r$ or AUC for Mann--Whitney, $\beta$ (with CI) for the linear model. Effect size $+$ CI $+$ $p$, together, is the honest report.

\paragraph{Aside --- entropy.} Where a distribution over $K$ categories (e.g.\ zonation states) is summarized, the Shannon entropy
\[
H = -\sum_{k=1}^{K} p_k \log p_k
\]
quantifies how spread-out the mass is: $H$ is maximal when all $p_k$ are equal and $0$ when one category holds all the mass. Concretely, with $K=4$ zonation states: a cell crisply assigned to one state ($p=(0.97,0.01,0.01,0.01)$) has $H\approx0.15$, while a smeared, ambiguous cell ($p=(0.25,0.25,0.25,0.25)$) hits the maximum $H=\log 4\approx1.39$. So high entropy flags a cell whose zonation identity is blurred --- exactly the per-cell signature of de-zonation we are hunting.

\bigskip

\noindent\fbox{\parbox{0.97\textwidth}{%
\textcolor{blue!45!black}{\textbf{One-page recap.}}\\[0.4em]
\textbf{H1 (trend, ordered):} Spearman $\rho$ for monotone trend; Jonckheere--Terpstra for the ordered-groups alternative; donor-level bootstrap for the CI; label-shuffle permutation as the negative control.\\[0.3em]
\textbf{H3 (association, controlled):} Mann--Whitney for two-group rank comparison; within-donor correlation then aggregate to defeat the stage confound; OLS $\plast \sim \dez + C(\text{stage}) + C(\text{donor})$ for the partial association.\\[0.3em]
\textbf{Everywhere:} donor $=$ unit of inference, BH-FDR across many tests, effect sizes beside $p$-values.}}

\end{document}
